% Independent GFDL mathematics, preceded by its own title/rights/history.
\noindent 四题的完整数学内容见下, 它们不依赖第1--17章的定理: 群是具有结合运算、单位元和双侧逆元的集合, 子群是同一运算下成群的子集. 原题号、历史页码和编者补明条件随题登记.

\begin{exercise}\label{ex:judson5}
MIT作业1原第3题; Judson第3章第5题, 两版印刷页52 (PDF第62页). 列出正方形的全部对称变换, 证明它们在复合下成群, 并写出乘法表. 四个顶点共有多少种置换? 每种置换是否都来自正方形对称变换? 此群记作 $D_4$.
\end{exercise}
\begin{solution}
按环向给顶点编号 $1,2,3,4$, 令 $r=(1234)$ 是四分之一周旋转, $s=(24)$ 是沿顶点1、3对角线的反射. 它们满足 $r^4=s^2=1$, $sr=r^{-1}s$. 全部对称变换恰为 $r^i$、$r^is$ ($0\le i<4$): 任一对称变换保持相邻关系, 顶点1的像有4种, 其相邻顶点2的像有2种; 顶点3的像是顶点1像的唯一对角点, 顶点4的像是顶点2像的唯一对角点, 因而全被确定. 三个不共线顶点的像唯一决定平面刚性变换. 上述8个变换恰实现这些选择, 故无遗漏.

刚性变换复合仍把正方形映到自身, 恒等变换在其中, 每个变换的逆也保持正方形; 函数复合满足结合律, 故形成8元群. 乘法约定行元素乘列元素, 右边先作用; 指数模4. 四类乘法分别为 $r^ir^j=r^{i+j}$, $r^i(r^js)=r^{i+j}s$, $(r^is)r^j=r^{i-j}s$, $(r^is)(r^js)=r^{i-j}$. 表如下:
\begin{center}\begin{tabular}{c|cccccccc}
$\cdot$&$1$&$r$&$r^2$&$r^3$&$s$&$rs$&$r^2s$&$r^3s$\\\hline
$1$&$1$&$r$&$r^2$&$r^3$&$s$&$rs$&$r^2s$&$r^3s$\\
$r$&$r$&$r^2$&$r^3$&$1$&$rs$&$r^2s$&$r^3s$&$s$\\
$r^2$&$r^2$&$r^3$&$1$&$r$&$r^2s$&$r^3s$&$s$&$rs$\\
$r^3$&$r^3$&$1$&$r$&$r^2$&$r^3s$&$s$&$rs$&$r^2s$\\
$s$&$s$&$r^3s$&$r^2s$&$rs$&$1$&$r^3$&$r^2$&$r$\\
$rs$&$rs$&$s$&$r^3s$&$r^2s$&$r$&$1$&$r^3$&$r^2$\\
$r^2s$&$r^2s$&$rs$&$s$&$r^3s$&$r^2$&$r$&$1$&$r^3$\\
$r^3s$&$r^3s$&$r^2s$&$rs$&$s$&$r^3$&$r^2$&$r$&$1$\\
\end{tabular}\end{center}
四顶点的全部置换共有 $4!=24$ 种, 其中仅上述8种来自对称变换. 例如 $(12)$ 固定3、4, 却把边 $(2,3)$ 的顶点变为对角线 $(1,3)$ 的顶点, 不保持距离, 因而不是正方形对称变换.
\end{solution}

\Needspace{10\baselineskip}
\begin{exercise}\label{ex:judson10}
MIT作业1原第4题; Judson第3章第10题, 两版印刷页52--53 (PDF第62--63页). 证明实矩阵集合
\[
\mathcal H=\left\{M(x,y,z)=\begin{pmatrix}1&x&y\\0&1&z\\0&0&1\end{pmatrix}:x,y,z\in\R\right\}
\]
在矩阵乘法下成群. 原题给出乘法公式, 此群称为海森堡群; 参数取实数是原章矩阵语境的显式化.
\end{exercise}
\begin{solution}
直接相乘得到 $M(x,y,z)M(x',y',z')=M(x+x',y+y'+xz',z+z')$, 故乘法封闭. 结合律也可直接核三元组: 三个因子的第二坐标无论怎样加括号均为 $y+y'+y''+xz'+xz''+x'z''$, 第一、三坐标分别为三个相应坐标之和. 单位元是 $M(0,0,0)=I_3$, 且
\[
M(x,y,z)^{-1}=M(-x,xz-y,-z).
\]
代入乘法式, 两个顺序的第二坐标分别为 $y+xz-y-xz=0$ 和 $xz-y+y-xz=0$, 其余坐标也为0. 因而每元有双侧逆, 群公理全部成立.
\end{solution}

\begin{exercise}\label{ex:judson48}
MIT作业1原第10题; Judson第3章第48题, 两版印刷页55 (PDF第65页). 群 $G$ 中元素 $a,b$ 满足 $a^4b=ba$ 且 $a^3=e$, 证明 $ab=ba$.
\end{exercise}
\begin{solution}
由 $a^3=e$ 得 $a^4=a^3a=a$. 将其代入第一个已知等式即得 $ab=ba$. 不需要先假定群交换.
\end{solution}

\begin{exercise}\label{ex:judson50}
MIT作业2原第14题 (挑战题); Judson第3章第50题, 两版印刷页55 (PDF第65页). 构造一个无限群, 使它的每个真子群都是有限群.
\end{exercise}
\begin{solution}
固定素数 $p$, 在复数乘法下取 $G=\bigcup_{n\ge0}G_n$, 其中 $G_n=\{z\in\C:z^{p^n}=1\}$. 这是准循环 $p$ 群. 令 $\zeta_n=\exp(2\pi i/p^n)$, 则 $G_n=\langle\zeta_n\rangle$, 恰有 $p^n$ 元, 且 $G_n\subseteq G_{n+1}$. 任意两个元素和其逆都落在某个共同的 $G_n$ 内, 因而 $G$ 是群; 它含任意大的有限子集 $G_n$, 故无限.

先说明循环群 $G_n$ 的子群都是某个 $G_k$ ($0\le k\le n$). 对非平凡子群 $H\le G_n$, 取最小正整数 $d$ 使 $\zeta_n^d\in H$. 用带余除法知 $d\mid p^n$, 且 $H=\langle\zeta_n^d\rangle$; 因为 $p$ 素, $d=p^{n-k}$, 正是 $G_k$. 平凡子群为 $G_0$.

现任取 $H\le G$. 若其元素阶有上界 $p^N$, 每元都是 $p^N$ 次单位根, 所以 $H\le G_N$, 是有限群. 若阶无界, 对每个 $n$ 可取 $h\in H$ 阶为 $p^m$, $m\ge n$. $h$ 位于某个 $G_j$ 中, 上述循环子群分类给出 $\langle h\rangle=G_m$, 从而 $G_n\subseteq H$. 于是所有 $G_n$ 都在 $H$ 中, 得 $H=G$. 因此真子群必有限, 且事实上恰为 $G_n$.
\end{solution}
